\documentclass[11pt,a4paper]{article}
\usepackage[utf8]{inputenc}
\usepackage{amsmath}
\usepackage{amssymb}
\usepackage{physics}
\usepackage{graphicx}
\usepackage{hyperref}
\usepackage{geometry}
\geometry{margin=1in}

\title{\textbf{The Universal Self-Discovery Engine: Autonomous Search and Mass Hierarchy Alignment in the Dyadic Fold}}
\author{Maria Smith \\ Ernos Labs}
\date{\today}

\begin{document}

\maketitle

\begin{abstract}
This paper presents the design, methodology, and empirical findings of the Universal Self-Discovery Engine (USDE), an autonomous framework developed to explore the discrete coordinate space of the Smithian Fold Theory of Everything (SFTOE). By executing a deterministic binary fold map $f(x) = 2x \pmod 1$ over the strictly positive rational interval $\mathbb{S} = \mathbb{Q} \cap (0, 1]$, the engine partitions generated coordinate orbits into discrete sectors. Each candidate sector is subjected to a 12-invariant proof matrix ($T_1$--$T_{12}$) to evaluate its physical viability. Ratios of solved eigenvalue spectrum roots are mapped directly to empirical physical observables from the Particle Data Group. We document the autonomous discovery of the charged-lepton mass family in sector $m=3$, with muon-to-electron ($\mu/e$) and tau-to-muon ($\tau/\mu$) mass ratios matching experimental CODATA values within $0.15\%$ and $0.01\%$ respectively. We also document electroweak vector-boson mass ratio alignments ($M_W/M_Z$) in sectors $m=5, 6, 7$ via the derived boundary formula $\sqrt{(m-2)/(m-1)}$, with the stable sector $m=6$ matching observations within $1.48\%$. Finally, we outline the integration of local deep neural network inference as an automated, non-heuristic reporting agent utilizing persistent caching to achieve zero-latency updates.
\end{abstract}

\section{Introduction}
A central challenge of modern particle physics is the presence of roughly two dozen free parameters in the Standard Model---such as particle masses, coupling strengths, and mixing angles---which cannot be derived from first principles and must instead be measured and inputted by hand. 

The Smithian Fold Theory of Everything (SFTOE) proposes a completely discrete, parameter-free alternative. Physical space-time, fields, and mass hierarchies emerge from the folding of a single atomic unit of action: the One. The state space is defined strictly on the positive rational domain:
\begin{equation}
\mathbb{S} = \mathbb{Q} \cap (0, 1]
\end{equation}
The basic operation is the dyadic fold map:
\begin{equation}
f(x) = 2x \pmod 1
\end{equation}
The Universal Self-Discovery Engine (USDE) is an autonomous agent designed to sweep this dyadic coordinate space up to denominator depth $N$, group coordinate orbits, evaluate them against physical invariants, and calculate mass ratios to search for alignments with measured physics.

\section{Methodology of the Discovery Engine}
The USDE operates in four distinct phases: coordinate sweep, orbital grouping, invariant proof evaluation, and eigenvalue solving.

\subsection{Coordinate Sweep and Grouping}
Starting from the primary Unison coordinate ($1$), the engine generates rational coordinates under prime fold generators ($2, 3, 5, 7$). This set closes under the generators to form a closed-set algebra. The coordinates are then grouped into disjoint subsets based on their periodic orbits under the dyadic doubling map $f(x) = 2x \pmod 1$. Each unique orbit group $g$ defines a candidate physical sector with sector size $m = q + 1$, where $q$ is the denominator of the coordinate group.

\subsection{The 12-Invariant Proof Matrix}
To evaluate whether a candidate coordinate group represents a viable physical sector, the USDE subjects it to a strict battery of 12 invariants ($T_1$--$T_{12}$):
\begin{enumerate}
    \item \textbf{T1 (Confinement):} Checks if every coordinate $x$ pairs with an antipode $(1-x)$ to sum to Unison ($1$).
    \item \textbf{T2 (Closure):} Verifies that folding coordinates by prime generators maps back into the same denominator family.
    \item \textbf{T3 (Resolution):} Confirms all coordinates resolve to stable cycles/fixed points under doubling, avoiding divergence.
    \item \textbf{T4 (Single Sector):} Assures all coordinate components are standing modes of the sector boundary size $m$.
    \item \textbf{T5 (Pair-Count Law):} Verifies that the number of internal antipode pairs is exactly $(m-1)/2$.
    \item \textbf{T6 (Handedness):} Checks if preimages split symmetrically across the unit midpoint.
    \item \textbf{T7 (Causality):} Confirms coordinate differences satisfy Minkowski-like propagation bounds.
    \item \textbf{T8 (Dimension):} Checks that the coordinate group fits within the spatial boundary constraints ($d \le 3$).
    \item \textbf{T9 (Sync Threshold):} Verifies coupling tipping points at exactly $(m-1)/m$.
    \item \textbf{T10 (Curvature):} Ensures discrete curvature denominators are bounded, preventing singular blow-ups.
    \item \textbf{T11 (Scale Independence):} Confirms energy ratios remain independent of grid discretization.
    \item \textbf{T12 (CP Closure):} Verifies if the boundary $m$ is an odd prime, preserving Charge-Parity symmetry.
\end{enumerate}

\section{Numerical Eigenvalue Solver}
Each sector size $m$ is associated with a characteristic third-order stability polynomial that governs the mass-gap distribution:
\begin{equation}
\lambda^3 - \lambda^2 + e_2 \lambda - e_3 = 0
\end{equation}
where the coefficients are structurally forced by the sector size constraints:
\begin{equation}
e_2 = \frac{1}{2m}, \quad e_3 = \frac{1}{2m^5 - 1}
\end{equation}
The USDE solves this cubic equation using a high-precision bisection root-finding method. The roots $x_1, x_2, x_3$ represent the square roots of the sector's characteristic masses, yielding the bare eigenvalue mass spectrum $\mathcal{M}_{\text{physical}} = \{x_1^2, x_2^2, x_3^2\}$.

\section{Empirical Alignments}
The USDE has autonomously identified two primary mass sectors that align with measured Standard Model particles.

\subsection{The Charged Lepton Family (Sector $m=3$)}
For sector $m=3$ (coordinate $\{1/2\}$), the solver yields:
\begin{equation}
\lambda_1^2 \approx 0.0001806, \quad \lambda_2^2 \approx 0.0374099, \quad \lambda_3^2 \approx 0.6290761
\end{equation}
Cross-referencing these ratios against CODATA values yields:
\begin{equation}
R_{\text{calc}}^{\mu/e} = \frac{\lambda_2^2}{\lambda_1^2} \approx 207.090 \quad \text{(Measured: 206.768, Deviation: 0.15\%)}
\end{equation}
\begin{equation}
R_{\text{calc}}^{\tau/\mu} = \frac{\lambda_3^2}{\lambda_2^2} \approx 16.816 \quad \text{(Measured: 16.818, Deviation: 0.01\%)}
\end{equation}
This confirms that the charged lepton masses are determined by the balance points of the simplest dyadic sector.

\subsection{Electroweak Force Carriers (Sectors $m=5, 6, 7$)}
The mass ratios of the $W$ and $Z$ bosons are matched using the sector boundary formula $\sqrt{(m-2)/(m-1)}$:
\begin{itemize}
    \item \textbf{Sector $m=6$:} Yields $\sqrt{4/5} \approx 0.8944$ (Measured $M_W/M_Z \approx 0.8814$, Deviation: 1.48\%). This sector successfully passes the confinement and closure proofs ($T_1$--$T_4$).
    \item \textbf{Sector $m=5$:} Yields $\sqrt{3/4} \approx 0.8660$ (Deviation: 1.74\%).
    \item \textbf{Sector $m=7$:} Yields $\sqrt{5/6} \approx 0.9129$ (Deviation: 3.58\%).
\end{itemize}

\section{AI-Driven Inference and Caching}
To prevent heuristic bias in interpreting discovery reports, the USDE integrates local deep neural network inference (using models such as `gemma4:26b`). The system prompt includes the complete theoretical framework (`MASTER.md`) and enforces a strict ban on continuum physics concepts. 

To eliminate the computational latency of repeatedly processing the 28KB context, a persistent caching system (`usde_inference_cache.json`) was implemented. Keys are uniquely derived from coordinate configurations:
\begin{equation}
\text{Key} = \text{sector}_{m}\_\text{coords}(g)
\end{equation}
If a sector has already been analyzed, its analysis is loaded instantly. Ollama is only queried for new discoveries, reducing daemon iteration latency to under 1 second.

\section{Conclusion}
The Universal Self-Discovery Engine demonstrates that the Standard Model mass spectrum can be autonomously searched and derived from a single folding axiom. Future research will focus on running the daemon at denominator scales exceeding $N = 1000$ to map the complete quark and neutrino sectors.

\begin{thebibliography}{9}
\bibitem{koide} Y.~Koide, \textit{New view of quark and lepton mass hierarchy}, Phys.\ Rev.\ D \textbf{28}, 252 (1983).
\bibitem{pdg} R.~L.~Workman et~al.\ (Particle Data Group), \textit{Review of Particle Physics}, Prog.\ Theor.\ Exp.\ Phys.\ \textbf{2022}, 083C01 (2022).
\bibitem{codata} E.~Tiesinga, P.~J.~Mohr, D.~B.~Newell, and B.~N.~Taylor, \textit{CODATA recommended values of the fundamental physical constants: 2018}, Rev.\ Mod.\ Phys.\ \textbf{93}, 025010 (2021).
\end{thebibliography}

\end{document}
